% =====================================================================
% JSAC2 v4
% On-device body digital twins for pose-aware mmWave links and live
% exposure readings.
%
% Author: Robin Wydaeghe
% Target: IEEE JSAC SI on Digital Twins for Wireless Networks (Q4 2026)
% Spine: ../JSAC_SPINE_PLAN_v2.md
% =====================================================================

\documentclass[journal,twocolumn,10pt]{IEEEtran}

\usepackage[utf8]{inputenc}
\usepackage[T1]{fontenc}
\usepackage{lmodern}
\usepackage{microtype}

\usepackage{amsmath,amssymb,amsthm,mathtools}
\usepackage{bm}
\usepackage{bbm}

\usepackage{graphicx}
\graphicspath{{figures/}{../code/outputs/}{../code/outputs/munich/}{../../theory/figures/}{../../JSAC/code/experiments/plaza_run/figures/}}
\usepackage{booktabs}
\usepackage{array}
\usepackage{multirow}
\usepackage{caption}
\usepackage{subcaption}

\usepackage{cite}
\usepackage{algorithm}
\usepackage{algpseudocode}
\usepackage{xcolor}
\usepackage[colorlinks=true,linkcolor=blue,citecolor=blue,urlcolor=blue]{hyperref}
\usepackage[capitalize,nameinlink]{cleveref}

\theoremstyle{plain}
\newtheorem{theorem}{Theorem}
\newtheorem{proposition}[theorem]{Proposition}
\newtheorem{corollary}[theorem]{Corollary}
\newtheorem{lemma}[theorem]{Lemma}
\theoremstyle{definition}
\newtheorem{definition}[theorem]{Definition}
\theoremstyle{remark}
\newtheorem{remark}[theorem]{Remark}

\crefname{proposition}{Approximation}{Approximations}
\Crefname{proposition}{Approximation}{Approximations}
\crefname{lemma}{Lemma}{Lemmas}
\Crefname{lemma}{Lemma}{Lemmas}
\crefname{theorem}{Theorem}{Theorems}
\Crefname{theorem}{Theorem}{Theorems}
\crefname{corollary}{Corollary}{Corollaries}
\Crefname{corollary}{Corollary}{Corollaries}
\crefname{remark}{Remark}{Remarks}
\Crefname{remark}{Remark}{Remarks}
\crefname{definition}{Definition}{Definitions}
\Crefname{definition}{Definition}{Definitions}

\newcommand{\khat}{\hat{\bm{k}}}
\newcommand{\nhat}{\hat{\bm{n}}}
\newcommand{\rhat}{\hat{\bm{r}}}
\newcommand{\etahat}{\hat{\bm{\eta}}}
\newcommand{\rr}{\mathbf{r}}
\newcommand{\xx}{\mathbf{x}}
\newcommand{\hh}{\bm{h}}
\newcommand{\EE}{\mathbf{E}}
\newcommand{\HH}{\mathbf{H}}
\newcommand{\Sinc}{S_{\mathrm{inc}}}
\newcommand{\Sab}{S_{\mathrm{ab}}}
\newcommand{\Pabs}{P_{\mathrm{ab}}}
\newcommand{\Aab}{A_{\mathrm{ab}}}
\newcommand{\Qmat}{\mathbf{Q}}
\newcommand{\Qabs}{\mathbf{Q}_{\mathrm{ab}}}
\newcommand{\Mmat}{\mathbf{M}}
\newcommand{\Jmat}{\mathbf{J}}
\newcommand{\Phimat}{\mathbf{\Phi}}
\newcommand{\Fn}{\mathbf{F}_n}
\newcommand{\psin}{\bm{\psi}_n}
\newcommand{\thetab}{\bm{\theta}}
\newcommand{\betab}{\bm{\beta}}
\newcommand{\gammab}{\bm{\gamma}}
\newcommand{\hBS}{\bm{h}_{\mathrm{BS}}}
\newcommand{\hbody}{\bm{h}_{\mathrm{b}}}
\newcommand{\Tepskin}{T_{0}}
\newcommand{\diff}{\mathrm{d}}
\newcommand{\Complex}{\mathbb{C}}
\newcommand{\Reals}{\mathbb{R}}
\newcommand{\ntilde}{\tilde{n}}
% (was: ReLU operator -- replaced by [\cdot]_+ notation throughout)
\DeclareMathOperator{\RE}{Re}
\DeclareMathOperator{\IM}{Im}
\DeclareMathOperator{\tr}{tr}
\DeclareMathOperator*{\argmin}{arg\,min}
\DeclareMathOperator*{\argmax}{arg\,max}

\title{On-Device Body Digital Twins for Pose-Aware mmWave Links\\
and Live Exposure Readings}

\author{Robin~Wydaeghe%
  \thanks{R. Wydaeghe is with Ghent University and imec, Ghent, Belgium
  (e-mail: robin.wydaeghe@ugent.be).}}

\markboth{IEEE Journal on Selected Areas in Communications, Month~2026}%
{Wydaeghe: On-device body digital twins}

\begin{document}

\maketitle

\begin{abstract}
Human bodies affect FR2 links in two measurable ways. A body in the
line of sight costs 25--40~dB of received power, while the same surface
absorbs part of the incident radio-frequency power. This paper proposes
an on-device digital twin of the user's body for both effects. The twin
uses an SMPL-X mesh with 66 pose degrees of freedom and a per-triangle
physical-optics kernel. One evaluation gives a live absorbed-power
density reading on the body surface and a pose-differentiable cascaded
channel $\hh(\thetab)=\hBS+\hbody(\thetab)$. The base station uses the
channel derivative to compute a small pose cue. The phone updates the
pose from inertial measurements, uplink pilots fit a low-rank
body-side correction, and the loop repeats. On 18 Sionna~RT Munich
macro-cell scenes (9~LOS, 9~NLOS), a naturalness-bounded pose move
gives 0.6--8.6~dB of cascaded-SINR gain. Mean gains are 1.7~dB on LOS
and 4.1~dB on NLOS. The strongest deep-shadow NLOS scene has 109~dB of
scene loss and gains 8.6~dB. Projected gradient ascent on a
32-dimensional pose latent reaches within 1~dB of the random-search
optimum in 20~steps on 16 of 18 scenes. At the canonical FR2 operating
point, the peak-local absorbed power density is 0.58\% of the ICNIRP
general-public limit.
\end{abstract}

\begin{IEEEkeywords}
Digital twin, closed-loop control, application-aware networks,
millimeter wave, MIMO, parametric body model, SMPL-X, physical optics,
Kirchhoff scattering, human-in-the-loop control, ICNIRP, absorbed
power density, electromagnetic-field exposure, RF-EMF epidemiology.
\end{IEEEkeywords}

\IEEEpeerreviewmaketitle


% =====================================================================
\section{Introduction}\label{sec:intro}
% =====================================================================

\IEEEPARstart{T}{he} FR2 millimeter-wave band (24--52.6~GHz) gives
5G and 6G systems the bandwidth needed for multi-Gbps peak
rates~\cite{general5G,TS38214}. Base-station arrays with tens to
hundreds of elements per panel deliver that bandwidth through
beamforming gain~\cite{HeathLozanoFA2018}.

The human body is part of the FR2 link budget. A single body in the
line of sight costs 25--40~dB of received
power~\cite{Maccartney2017blockage,MacCartney2017}. Panel array gain
does not recover that loss once a pedestrian blocks the path. The body
also absorbs radio-frequency power. ICNIRP~2020 specifies absorbed
power density above 6~GHz~\cite{ICNIRP2020}, and local field caps have
affected millimeter-wave rollout in the Brussels Capital
Region~\cite{BrusselsArrete2007,BrusselsArrete2024}, several Italian
regions~\cite{ItalyDL}, and Switzerland under federal
limits~\cite{SwissNISV2000}.

Prior work treats the two effects separately. Exposure-aware
beamforming writes absorbed power as the Hermitian quadratic form
$\Pabs=\xx^{H}\Qabs\xx$, where $\Qabs$ depends on the body and pose
and $\xx$ is the precoder~\cite{Hochwald2014,Ying2015,Ying2017}.
Reconfigurable intelligent surface (RIS) models describe the cascaded
channel inserted by a passive scatterer between the base station and
the user equipment~\cite{DiRenzo2020,Wu2019}. A human body is also a
passive lossy scatterer. Its controllable state is pose, and the user
is the actuator.

That actuator receives little useful information today. The user does
not see that their body blocks the link. The phone does not show where
radio-frequency power is absorbed on the body surface. The network
does not return a pose cue that improves the link while keeping the
movement small. A body digital twin can make these quantities local,
pose-resolved, and usable in the radio loop.

Three gaps limit current models of pose-aware FR2 links:
(1)~No current body digital twin gives users, regulators, or
epidemiologists an instantaneous exposure reading on the body surface.
(2)~No current channel model treats the body-mediated coefficient as a
single-receiver render anchored at the phone, which is the geometry
available on device. (3)~The closed loop between on-device pose
telemetry, twin update, and base-station precoder adaptation has not
been studied end to end on a ray-traced millimeter-wave scene.

This work addresses these gaps with an on-device digital twin of the
user's body inside the per-slot radio control loop. The twin is a
parametric body mesh~\cite{SMPLX2019}. The phone, the base station,
and a ray-traced scene~\cite{Hoydis2023Sionna} carry the propagation.
The same per-triangle physical-optics kernel produces two outputs from
one evaluation: a live absorbed-power-density map on the body surface
and a pose-differentiable cascaded channel. The base station computes
a small pose cue from the channel derivative. The user moves, the
inertial measurement unit (IMU) and uplink pilots update the pose and
channel, a low-rank body-side coefficient absorbs the residual, and
the loop repeats.

To the best of the authors' knowledge, the following are reported here
for the first time.
\begin{enumerate}
\item A body digital twin that closes a per-slot control loop with a
ray-traced FR2 base station, end-to-end pose differentiability of the
rate, and a human-in-the-loop law that converts the base-station
gradient into a sub-degree pose cue
(\Cref{sec:arch,sec:control}).
\item A receiver-anchored Kirchhoff render of the body-mediated
channel $\hbody(\thetab)$ that runs at consumer-grade compute budgets
on a single phone (\Cref{sec:kirchhoff}).
\item A whole-body absorbed-power-density read-out on the same per-pose
budget, four orders of magnitude faster than the IT'IS
Virtual-Population FDTD reference~\cite{Christ2010VF} at comparable
accuracy on biological tissue, across the full FR2
band (\Cref{sec:saridx}).
\item An empirical study on 18~Sionna~RT scenes of the Munich
macro-cell (9~LOS, 9~NLOS) that reports existence and reachability of
naturalness-bounded pose gains and a failure-mode taxonomy
(\Cref{sec:numerical}).
\end{enumerate}

\begin{table}[!t]
\caption{Two semantic outputs from one per-triangle render.}
\label{tab:two-products}
\centering\small
\begin{tabular}{@{}p{0.30\columnwidth}p{0.27\columnwidth}p{0.30\columnwidth}@{}}
\toprule
& \textbf{Live exposure reading} & \textbf{Pose-driven capacity} \\
\midrule
Computes & $\Sab(\rr)$ on body surface & $\nabla_{\thetab} R(\thetab,\xx)$ \\
Output & dose number, displayed locally & pose move, suggested via DL \\
Where it lives & on the body surface & at the phone antenna \\
Privacy delta & reading local, opt-in aggregation & gradient local, only $\thetab^{*}$ accept/reject leaves the device \\
Loop role & informational & rate maximization under comfort budget \\
\bottomrule
\end{tabular}
\end{table}


% =====================================================================
\section{System model and body-twin loop}\label{sec:arch}
% =====================================================================

\subsection{Geometry and channel}

The body-twin loop is shown in~\Cref{fig:dtn-arch}. A base
station serves a single user terminal in the FR2 band, the user
holds a phone, and the pose-mediated body channel is rendered on
the device. The remainder of this section fixes the geometry, the
channel split, and the cadence at which the loop runs.

The base station is a rectangular array of $M$ elements with
half-wavelength spacing, operating at carrier $f_c$ in the FR2 band
($f_c = 28$~GHz throughout, except where indicated). The user is a
single served terminal at world position $\rr_p \in \Reals^{3}$
holding a phone with one receive antenna, isotropic in pattern for
simplicity.\footnote{Polarization- and pattern-aware extensions
follow via the projection $\bm{g}_{\mathrm{UE}}^{H}$ introduced
in~\Cref{sec:kirchhoff} and are suppressed in the main derivation.}
The cascaded narrowband channel vector at the phone, with $M$
complex entries (one per BS element), is
\begin{equation}\label{eq:cascaded}
\hh(\thetab,\rr_p) \;=\; \hBS(\rr_p) \;+\; \hbody(\thetab,\rr_p) \,,
\end{equation}
with $\hh,\hBS,\hbody \in \Complex^{M}$. The first term $\hBS$
collects every BS-to-phone propagation path that does not interact
with the body, that is the direct LOS when present and all
scene-side multipath via building reflections and edge diffraction.
The second term $\hbody$ is the body-mediated coefficient computed
by the on-device twin. The split is exact by Maxwell linearity.

The base station transmits the signal $\xx s$, where the precoder
$\xx \in \Complex^{M}$ satisfies $\|\xx\|^{2} \leq P_{\mathrm{tx}}$
and $s$ is a unit-power data symbol. The received signal at the
phone is
\begin{equation}\label{eq:rx}
y \;=\; \hh^{H} \xx \, s + n \,, \qquad n \sim \mathcal{CN}(0, N_{0}) \,,
\end{equation}
and the achievable rate per Hz is $R = \log_{2}(1 + \mathrm{SINR})$
with SINR defined in~\Cref{sec:control}. Rates are reported in
bps/Hz throughout, with any modulation-coding-scheme cap left
implicit at the implementation layer.

\subsection{Body twin and pose telemetry}

The user's body is an SMPL-X parametric
mesh~\cite{SMPLX2019} with $22$ body joints, $10\,475$ vertices,
and $20\,908$ triangles. The pose state is the axis-angle vector
$\thetab \in \Reals^{22 \times 3}$. Body-shape parameters
$\bm{\beta}$, distinct from the BS-side calibration coefficient
$\betab$ of~\Cref{sec:closed-loop}, are fixed user attributes
calibrated once at enrolment. The forward map $\thetab \mapsto
\Sigma(\thetab)$ is differentiable through linear blend skinning,
which makes the gradient $\nabla_{\thetab} R$
of~\Cref{sec:control} a tractable object.

The BS holds a separate environmental twin: a ray-traced
description of the scene~\cite{Hoydis2023Sionna} from which a
per-element path dictionary $\Jmat \in \Complex^{M \times P}$ is
extracted, with $P$ the number of dominant multipath components
(MPCs) at the served-user location. Each column of $\Jmat$ is one
MPC. The dictionary $\Jmat$ is the only per-user object the BS
needs from the scene, refreshed at the user-mobility timescale,
and feeds both the body-side decomposition of~\Cref{sec:kirchhoff}
and the precoder design of~\Cref{sec:control}.

Pose telemetry is supplied by an in-pocket consumer-grade IMU
running a Madgwick-style attitude filter~\cite{Madgwick2011}, with
steady-state per-joint attitude RMS error
$\sigma_{\mathrm{joint}} \approx 4^{\circ}$ at the canonical
operating point. The IMU error model and the
$\sigma_{\mathrm{joint}}$ sweep are deferred
to~\Cref{sec:closed-loop} and the supplementary information. Pose
enters the twin via the IMU stream and is corrected over the air
by the closed-loop calibration of~\Cref{sec:closed-loop}.

\subsection{Closed loop}

\begin{figure*}[!t]
\centering
\includegraphics[width=0.92\textwidth]{figures/fig_dtn_architecture.pdf}
\caption{Network-centric body-twin loop. The base station
maintains the per-element path dictionary $\Jmat$ and the precoder
$\xx$ solver. The on-device body twin holds the SMPL-X mesh
$\Sigma(\thetab)$, runs the UE-anchored Kirchhoff render to obtain
$\hbody(\thetab,\rr_p)$, fits the body-side calibration
coefficient $\gammab$ over a $K_{99}$-mode SVD basis from
uplink-pilot residuals, and produces both the rate gradient
$\nabla_{\thetab} R$ and the absorbed-power-density field
$\Sab(\rr)$. The user is the actuator of $\thetab$ and the
substrate of the IMU. Solid arrows are forward (downlink), dashed
arrows are return (uplink). Two semantic outputs land below their
producers, a network-side rate adaptation and a user-side live
exposure reading.}
\label{fig:dtn-arch}
\end{figure*}

At each twin tick, the loop runs five steps.
\begin{enumerate}
\item The base station refreshes the path dictionary $\Jmat$ at
$\sim\!1$~Hz, and per-path amplitude calibration $\betab$ runs
over uplink pilots at the same cadence.
\item The on-device twin reads IMU pose $\hat\thetab$, runs the
SMPL-X forward map to obtain the mesh, evaluates the UE-anchored
Kirchhoff render of~\Cref{sec:kirchhoff} to obtain $\hbody$, and
combines it with $\hBS$ to produce a cascaded-channel estimate.
\item The BS observes the served user's actual channel via uplink
pilots, computes the residual against the twin's prediction, and
fits the body-side coefficient $\gammab$ over a $K$-mode SVD basis
of the body-to-phone Kirchhoff operator
(\Cref{sec:closed-loop}).
\item The BS computes a pose-direction suggestion
$\nabla_{\thetab} R(\thetab) / \|\nabla_{\thetab} R\|$ truncated to
a comfortable magnitude and emits a low-bandwidth packet over the
downlink control channel ($\sim\! 6 M r$ float scalars per
suggestion at low rank $r$, see~\Cref{sec:bandwidth}).
\item The user, optionally, actuates a small pose move toward the
suggestion. The loop closes at the next twin tick.
\end{enumerate}
The on-device twin separately produces the absorbed-power-density
reading $\Sab(\rr;\thetab,\xx)$ for display to the user and (opt-in)
upload to a long-running RF-EMF cohort study. This branch is
informational and does not enter the rate-optimization loop.

\subsection{Positioning against the digital-twin literature}

Digital-twin work in wireless networks usually twins the network
infrastructure, an industrial process, or an autonomous
agent~\cite{Chen2024DTN,Saad2025DTN,Hashash2025DTN}. Here the twin
is the human body in the cell. The live exposure reading and the
pose-gradient action are two outputs of one physical model. Both
directions of the loop use the same per-triangle Kirchhoff render
of~\Cref{sec:kirchhoff}.


% =====================================================================
\section{Body physics: per-triangle Fresnel and PO scattering}\label{sec:physics}
% =====================================================================

This section gives the per-triangle scattering primitive that
both outputs of~\Cref{tab:two-products} use. The model is a
single-bounce Fresnel physical-optics model on a lossy half-space at
mmWave. Polarization-aware Stokes-vector dosimetry, curvature
corrections, diffraction smoothing, and inter-body radiosity
extensions are deferred to the supplementary information.

\subsection{Fresnel coefficients at a lossy half-space}

\begin{figure}[!t]
\centering
\includegraphics[width=0.95\columnwidth]{incidence_plane_fig.pdf}
\caption{Per-triangle incidence-plane geometry on the body surface
$\Sigma$. The TE basis $\hat{e}_{s,n}$ (green) is orthogonal to the
incidence plane. The incident TM basis $\hat{e}_{p,n}$ (orange) lies
in the plane with large normal component $\sin\theta_n$. The
refracted ray $\hat{\bm{k}}'_n$ (dashed grey) enters at the smaller
angle $\theta_t$ because of the high refractive index, and the
refracted TM direction $\hat{e}'_{p,n}$ (blue) is nearly tangential
($\sin\theta_t = O(1/|\ntilde|)$). Approximation~1
(\Cref{prop:approx1}) replaces $\hat{e}'_{p,n}$ by $\hat{e}_{p,n}$.
The cross-term error is $O(1/|\ntilde|^2)$ because the small normal
component dominates the difference.}
\label{fig:incidence}
\end{figure}

Consider a planar interface between vacuum and a half-space of
biological tissue with complex refractive index $\ntilde =
\sqrt{\varepsilon_r - j\sigma/(\omega\varepsilon_0)}$. For skin at
$28~\mathrm{GHz}$, $\ntilde \approx 4.49 - 1.79j$ and
$|\ntilde| \approx 4.83$. A plane wave incident at angle $\theta_n$
from the surface normal $\nhat$ admits the standard Fresnel TE and
TM amplitude coefficients $t_{s,n}$ and $t_{p,n}$, derived in
the supplementary information.

The incident wave $\psin = \psi_{s,n}\hat{e}_{s,n} +
\psi_{p,n}\hat{e}_{p,n}$ excites the transmitted field
$\EE^{\mathrm{trans}}_n$ inside the tissue, and the absorbed power
density at the surface, integrated over the skin depth, is
\begin{equation}\label{eq:Sab-single}
  \Sab^{(n)}(\rr) = \frac{|x_{j(n)}|^2}{Z_0}\bigl(T_{s,n}|\psi_{s,n}|^2
  + T_{p,n}|\psi_{p,n}|^2\bigr)[\mu_n]_+\,,
\end{equation}
with $T_{s,n} = |t_{s,n}|^2\,\RE(\xi_n)/\mu_n$ the TE power
transmittance, $T_{p,n}$ the TM analog, $\xi_n = \sqrt{\ntilde^2-1+\mu_n^2}$,
$Z_0 \approx 377\,\Omega$ the free-space impedance,
$\mu_n = -\hat{\bm{k}}_n\cdot\nhat$ the incidence cosine, and
$[\mu_n]_+$ the back-face cull.

\subsection{Pseudo-Brewster collapse}

For materials with $|\ntilde| \geq 2 + \sqrt{3} \approx 3.73$, the
TE and TM power transmittances pinch together at all incidence
angles below the critical angle. We refer to this effect as the
\emph{pseudo-Brewster collapse}. Skin at mmWave satisfies the
condition. At 28~GHz, both transmittances lie within 5\% of a single
scalar
$\Tepskin \approx 0.54$ across
$\theta_n \in [0^{\circ}, 80^{\circ}]$. Substituting into
\Cref{eq:Sab-single} collapses the polarization structure to
$\Sab^{(n)} = \Sinc^{(n)}\,\Tepskin\,[\mu_n]_+$ with
$\Sinc^{(n)} = |x_{j(n)}|^2\,|\psin|^2 / Z_0$, the geometric
absorption law on a lossy half-space~\cite{BornWolf}. The
polarization-aware extension is retained only for error bounds.

\subsection{Body as a discretized PO scatterer}

The user's body is the SMPL-X mesh of~\Cref{sec:arch}: a closed
surface $\Sigma(\thetab)$ tessellated into $T = 20\,908$ triangles.
Triangle $t$ at pose $\thetab$ has centroid $\bm{c}_t(\thetab)$,
outward unit normal $\nhat_t(\thetab)$, and area $S_t(\thetab)$.
The visibility factor for incident direction $\hat{\bm{k}}_n$ is
\begin{equation}\label{eq:visibility}
  V_n(t) = [-\hat{\bm{k}}_n \cdot \nhat_t]_+\,,
\end{equation}
the standard back-face cull. A strict ray-mesh visibility test is
applied only when fingertip-scale features matter
(SI~\S{}S8).

The base-station path dictionary $\Jmat \in \Complex^{M \times N}$
encodes $N$ Sionna~RT-traced propagation paths from the BS to the
body region. Each path $n$ carries a per-element steering response
$\bm{a}_n \in \Complex^{M}$, an arrival direction $\hat{\bm{k}}_n$,
a polarization $\psin$, and a per-path amplitude $\beta_n$
(calibrated against UL pilots, see SI~\S{}S5). For triangle $t$
and path $n$, the incident wave at the body surface is
\begin{equation}\label{eq:incident-wave}
  \psin\,\beta_n\,\bm{a}_n^{H}\xx\,
  e^{-jk_0 \hat{\bm{k}}_n \cdot \bm{c}_t}\,.
\end{equation}

\subsection{Exposure operator and two enabling approximations}

The whole-body absorbed power
$\Pabs(\thetab,\xx) = \int_{\Sigma}\Sab\,\mathrm{d}A$ is a Hermitian
quadratic form,
\begin{equation}\label{eq:Qabs}
  \Pabs(\thetab,\xx) = \xx^{H}\,\Qabs(\thetab)\,\xx\,,
\end{equation}
under two approximations whose error budgets are tight on
biological tissue at mmWave.

\begin{proposition}[TM polarization direction; Approximation~1]\label{prop:approx1}
The exact transmitted TM field points along the refracted direction
$\hat{e}'_{p,n}$ inside the tissue. Replace $\hat{e}'_{p,n}$ by the
incident direction $\hat{e}_{p,n}$ (\Cref{fig:incidence}). Self-terms
are unaffected, since $|\hat{e}_{p,n}|^2 = |\hat{e}'_{p,n}|^2 = 1$.
The error on TM--TM cross-terms is $O(1/|\ntilde|^2) \lesssim 4\%$
for skin at 28~GHz, and shrinks at higher frequencies and
on tissues with larger $|\ntilde|$.
\end{proposition}

\begin{proposition}[Universal depth coupling; Approximation~2]\label{prop:approx2}
The depth-integrated cross-coupling between paths $n$ and $n'$ is
\begin{equation}\label{eq:Lambda}
  \Lambda_{nn'} = \frac{1}{\alpha_n + \alpha_{n'}
  - j(\beta_{n'} - \beta_n)}\,,
\end{equation}
with $\alpha_n, \beta_n$ the per-path attenuation and phase rates
inside the tissue. Define the normalized coupling
$\Gamma_{nn'} = \Lambda_{nn'} / \sqrt{\Lambda_{nn}\Lambda_{n'n'}}$,
and approximate $\Gamma_{nn'} \approx 1$. For skin at
28~GHz over all incidence-angle pairs
$(\theta_n,\theta_{n'}) \in [0^{\circ},85^{\circ}]^2$, the error
satisfies $|1 - \Gamma_{nn'}| \leq 0.44\%$ with mean $0.17\%$.
The high refractive index compresses refraction angles to within
$\sim\!10^{\circ}$ of normal, so $\alpha_n,\beta_n$ are nearly
angle-independent.
\end{proposition}

Proofs are given in the supplementary information. The combined
effect of~\Cref{prop:approx1,prop:approx2} is to factor the local
absorbed power as a squared norm,
\begin{equation}\label{eq:Sab-norm}
  \Sab(\rr) = \bigl\|\tilde{\bm{G}}(\rr)\,\xx\bigr\|^2\,,
\end{equation}
with $\tilde{\bm{G}}(\rr) \in \Complex^{3 \times M}$ the
\emph{exposure channel matrix}. Its $j$-th column reads
\begin{equation}\label{eq:exposure-channel}
  \tilde{\bm{g}}_j(\rr) = \!\!\sum_{n:\,j(n)=j}\!\!
  \sqrt{\sigma/4\alpha_n}\,
  \Fn(\rr)\,\psin\,e^{-jk_0\hat{\bm{k}}_n \cdot \rr}\,,
\end{equation}
with $\Fn$ the per-path Fresnel transmission operator (rank-2,
polarization-aware, formally defined in the supplementary
information). Integration over the body surface gives
\begin{equation}\label{eq:Q-form}
  \Qabs(\thetab) = \!\!\int_{\Sigma(\thetab)}\!\!
  \tilde{\bm{G}}(\rr)^{H}\,\tilde{\bm{G}}(\rr)\,\mathrm{d}A
  = \Jmat^{H}\,\Mmat(\thetab)\,\Jmat\,,
\end{equation}
where $\Mmat(\thetab) \in \Complex^{N\times N}$ is the
body-and-scene Gram matrix and $\Jmat$ the array-hardware path
dictionary. The factorization~\eqref{eq:Q-form} separates the
per-slot precoder cost from the per-pose body refresh, the
structural property that makes BS-side exposure-constrained
beamforming cheap~\cite{Hochwald2014,Ying2015}. The same
factorization carries over to the exposure-reading branch.


% =====================================================================
\section{UE-anchored Kirchhoff render}\label{sec:kirchhoff}

The UE-anchored Kirchhoff render is illustrated in
\Cref{fig:kirchhoff-schematic}. Path $n$ from the BS path dictionary
illuminates triangle $t$ on the body surface. The local Fresnel
reflection re-radiates a field whose contribution at the phone
$\rr_p$ is the per-point integrand $\mathcal{K}_n(\rr;\rr_p)$. The
body-mediated channel coefficient $\hbody(\thetab,\rr_p)$ is the
integral of this kernel over the visible surface, summed over BS
paths.

\begin{figure}[!t]
\centering
\includegraphics[width=\columnwidth]{figures/fig_kirchhoff_schematic.pdf}
\caption{Per-triangle UE-anchored Kirchhoff render. The base station
illuminates triangle $t$ via path $n$ with arrival direction
$\khat_n$. The local Fresnel reflection re-radiates toward the
phone at $\rr_p$ along the per-point arrival direction
$\etahat(\rr;\rr_p)$. The body-mediated channel coefficient
$\hbody(\thetab,\rr_p)$ is the integral of the per-point kernel
$\mathcal{K}_n$ over the visible surface, summed over BS paths.}
\label{fig:kirchhoff-schematic}
\end{figure}

% =====================================================================

The render has two steps. At each surface point on
$\Sigma(\thetab)$, the local Fresnel operator reflects the incident
BS multipath into a re-radiated field. The free-space Green's
function then propagates that field to the phone position $\rr_p$,
which gives a Kirchhoff radiation integral.

\subsection{PO surface current and visibility}

For path $n$ incident on triangle $t$, the local reflected current
density takes the standard physical-optics form $\bm{J}^{(n)}_t =
2\,(\nhat_t \times \HH^{(n)}_t)\,V_{nt}$, with $\HH^{(n)}_t$ the
local incident magnetic field (derived from $\psin$ via the
impedance $Z_0$) and $V_{nt} \in \{0,1\}$ the per-triangle, per-path
visibility flag. Only triangles with $V_{nt}=1$ contribute to the
Kirchhoff integral~\eqref{eq:kirchhoff-integral}. The back-face
cull $\nhat_t \cdot \khat_n < 0$ suffices for the phone-on-body
geometries treated here. The full ray-mesh self-shadow test is
reserved for fingertip features and is detailed in the
supplementary information.

\subsection{Propagation to the phone}

The free-space Green's function from a surface point $\rr$ to the
phone at $\rr_p$ is
\begin{equation}
G(\rr,\rr_p) = \frac{\exp(jk_0\,\|\rr-\rr_p\|)}{4\pi\,\|\rr-\rr_p\|} \,,
\end{equation}
and the per-point arrival direction at the phone is
$\etahat(\rr;\rr_p) = (\rr_p - \rr)/\|\rr_p - \rr\|$. The
body-mediated channel coefficient then reads
\begin{equation}\label{eq:kirchhoff-integral}
\hbody(\thetab,\rr_p) = \sum_{n=1}^{N} \bm{a}_n
\int_{\Sigma_n^{+}(\thetab)} \mathcal{K}_n(\rr;\rr_p)\,\mathrm{d}A \,,
\end{equation}
with per-point integrand
\begin{equation}\label{eq:kernel}
\mathcal{K}_n(\rr;\rr_p) = K\,\bm{g}_{\mathrm{UE}}(\etahat)^{H}\psin
\,G(\rr,\rr_p)\,\exp(-jk_0\,\khat_n\cdot\rr) \,.
\end{equation}
Here $K = jk_0 r_0$ is the pseudo-Brewster reflection amplitude with
$r_0 = 1 - \Tepskin \approx 0.46$, and $\bm{g}_{\mathrm{UE}}(\etahat)$
is the phone's antenna pattern in the arrival direction (isotropic
by assumption, $|\bm{g}_{\mathrm{UE}}|^2 \equiv 1$). The visible
region $\Sigma_n^{+}$ collects the surface points that satisfy the
back-face cull for path $n$ and lie in line of sight of $\rr_p$.

\subsection{Discretization and Kirchhoff operator}

Replacing each integral over $\Sigma_n^{+}$ by a sum over the
$T_{\mathrm{vis}}$ visible triangles yields the matrix form
\begin{equation}\label{eq:Lambda-op}
\hbody(\thetab,\rr_p) = \Lambda_{\mathrm{KH}}(\thetab,\rr_p)\,
\bm{1}_{T_{\mathrm{vis}}} \,,
\end{equation}
with $\Lambda_{\mathrm{KH}} \in \Complex^{M \times T_{\mathrm{vis}}}$
the \emph{body-to-phone Kirchhoff operator}. Its $(m,t)$ entry is
$\Lambda_{\mathrm{KH}}[m,t] = \sum_{n \to t} [\bm{a}_n]_m \,
\beta_n\, S_t\, \mathcal{K}_n(\bm{c}_t;\rr_p)$, where $\bm{c}_t$
is the centroid and $S_t$ the area of triangle $t$, and the sum
runs over BS paths landing on $t$. The operator
$\Lambda_{\mathrm{KH}}$ is the central object of the body-channel
prediction. Its singular spectrum, analyzed in
\Cref{sec:closed-loop}, governs calibration identifiability.

\begin{remark}[Translation phasor identity]
For a small body displacement $\bm{t}$ (e.g., a walking step), the
operator transforms as $\Lambda_{\mathrm{KH}}(\rr_0 + \bm{t}) =
\Phimat(\bm{t})^{H}\Lambda_{\mathrm{KH}}(\rr_0)$, with $\Phimat(\bm{t})$
a diagonal phasor matrix on the BS-element axis. This is the analog
of the standard array-translation phase identity. The per-slot
refresh under translation then costs only one diagonal multiply.
\end{remark}

The render runs at $\sim\!100\,$ms per pose on a single CPU core
at $T = 20\,908$, $T_{\mathrm{vis}} \sim 10\,000$, and $M = 64$ on
the SMPL-X body of \Cref{sec:arch}.


% =====================================================================
\section{Closed-loop calibration: UL pilots and \texorpdfstring{$\gammab$}{gamma}}\label{sec:closed-loop}
% =====================================================================

The Kirchhoff render of~\Cref{sec:kirchhoff} is a \emph{prediction}
of $\hbody$ from the on-device twin. The over-the-air channel
differs from this prediction for three reasons. First, the SMPL-X
mesh is not the user's exact anatomy. Second, the dielectric
assumption $\Tepskin\approx 0.54$ is a population mean. Third, the
IMU pose is not the ground-truth pose. The calibration step closes
this gap by regressing the measured channel on the predicted one,
with a refit per slot or per second from the uplink pilot stream.

\subsection{Residual decomposition}

Let $\hh^{\mathrm{meas}}$ be the channel observed at the BS from
uplink pilots and $\hh^{\mathrm{pred}}(\thetab) =
\hBS + \hbody^{\mathrm{pred}}(\thetab)$ the cascaded channel
predicted by the twin. The residual $\Delta = \hh^{\mathrm{meas}} -
\hh^{\mathrm{pred}}$ splits into BS-side and body-side components,
\begin{equation}\label{eq:residual}
\Delta = \underbrace{\Phimat(\rr_p)\Jmat\,\delta\betab}_{\text{BS-side}}
+ \underbrace{\Lambda_{\mathrm{KH}}(\thetab,\rr_p)\,\delta\bm{q}}_{\text{body-side}}
+ \bm{n}_{\mathrm{p}}\,,
\end{equation}
where $\delta\betab$ collects the per-path amplitude corrections
(fitted from uplink pilots, SI~\S{}S5), $\delta\bm{q}$ is a
per-triangle correction on the body-mediated current, and
$\bm{n}_{\mathrm{p}}$ is the pilot noise. Both sides are
identifiable from the same uplink stream when the column spaces of
$\Phimat\Jmat$ and $\Lambda_{\mathrm{KH}}$ are linearly independent.

\subsection{Body-side parameterization: \texorpdfstring{$K$}{K}-mode SVD basis}

The body-side correction $\delta\bm{q}$ has $T_{\mathrm{vis}} \sim
10^{4}$ entries, far more than the $M = 64$ pilot equations admit.
We project onto a $K$-dimensional basis $\bm{V}_K \in
\Complex^{T_{\mathrm{vis}}\times K}$ and write $\delta\bm{q}
\approx \bm{V}_K \gammab$. The coefficient $\gammab \in
\Complex^{K}$ then follows from a ridge least-squares fit against
the residual. We take $\bm{V}_K$ as the leading $K$ right singular
vectors of $\Lambda_{\mathrm{KH}}$, with $K = K_{99}$ the smallest
$K$ that captures $99\%$ of the spectral mass of the operator.
The operator-spectrum sweep of SI~\S{}S7 reports $K_{99}$ ranging
from $1$ to $4$ across the BS geometries considered, with the
exact value depending on UE--BS distance and visibility.

\begin{proposition}[Identifiability]\label{prop:ident}
Let $\bm{U}_K$ be the leading-$K$ left singular vectors of
$\Lambda_{\mathrm{KH}}$, with corresponding singular values $\sigma_k
> 0$. The body-side ridge least squares with basis $\bm{V}_K$ equal
to the corresponding right singular vectors is exactly
identifiable when the residual lies in $\mathrm{span}(\bm{U}_K)$.
The fit is non-trivial as long as $K \leq M$ and the pilot SNR
exceeds $K \sigma_K^{-2}$.
\end{proposition}

The proof is a standard truncated SVD argument and is given in
the supplementary information.

\subsection{Calibration robustness to IMU drift}

The body-side $\gammab$ absorbs pose error (residual IMU drift
after attitude filtering) and dielectric mismatch in the same
regression. To first order, both error sources produce residuals
in the column space of $\Lambda_{\mathrm{KH}}$. The same $K$-mode
SVD basis that captures the dominant body-to-phone modes therefore
captures the dominant residual directions, and the calibration
step recovers the predicted channel from a noisy pose input by
projecting the residual onto the leading modes of the operator.
This projection makes the closed loop robust to consumer-grade IMU
drift. The IMU-noise sensitivity sweep of SI~\S{}S6 quantifies the
rate cost over $\sigma_{\mathrm{joint}} \in
\{0,2,4,8,16\}^{\circ}$.


% =====================================================================
\section{Pose-differentiable rate and human-in-the-loop control}\label{sec:control}
% =====================================================================

\subsection{Cascaded SINR and rate}

The $M$-element BS applies maximum-ratio transmission
$\xx = \sqrt{P_{\mathrm{tx}}}\,\hh / \|\hh\|$. The post-precoding
SINR at the phone reads
$\mathrm{SINR}(\thetab) = P_{\mathrm{tx}} \|\hh(\thetab)\|^{2} / N_0$,
so the achievable rate per Hz is
\begin{equation}\label{eq:rate}
R(\thetab) = \log_2\bigl(1 + \mathrm{SINR}(\thetab)\bigr)\,.
\end{equation}
The cascaded structure of~\eqref{eq:cascaded} carries the full
pose-dependence. The gradient $\nabla_{\thetab} R$ is non-zero
wherever $\mathrm{SINR}>0$, with magnitude
$\partial R / \partial \mathrm{SINR} = (\ln 2)^{-1}/(1 + \mathrm{SINR})$.

\subsection{Pose-space restriction and latent gradient}\label{sec:vposer-latent}

The axis-angle vector $\thetab \in \Reals^{66}$ is over-parameterized.
Most of $\Reals^{66}$ corresponds to anatomically impossible
configurations such as joints outside their range of motion or
incompatible torso-limb couplings. Optimization on $\Reals^{66}$
would routinely propose contortionist poses that the user cannot
actuate. We restrict the optimization to the manifold of plausible
body poses parameterized by the VPoser variational
prior~\cite{Pavlakos2019VPoser}. VPoser is a variational autoencoder
trained on the Archive of Motion Capture as Surface Shapes
(AMASS) corpus~\cite{AMASS2019}. Its
decoder $\bm{G}: \Reals^{32} \to \Reals^{63}$ maps a 32-dimensional
latent $\bm{z}$ to the 21 SMPL-X body-pose joints (excluding the
global orientation and the hands), and the prior
$p(\bm{z}) = \mathcal{N}(\bm{0}, \bm{I})$ assigns high probability
to natural poses. Every latent $\bm{z}$ decodes to an anatomically
valid configuration.

The pose-differentiable rate becomes a function of the latent
through $\thetab(\bm{z}) = \bm{G}(\bm{z})$,
$\hh(\bm{z}) = \hBS + \hbody(\thetab(\bm{z}))$, and
$R(\bm{z}) = \log_2(1 + P_{\mathrm{tx}}\|\hh(\bm{z})\|^{2}/N_0)$.
The differentiable map is
\begin{equation}\label{eq:autodiff-composition}
\bm{z} \xrightarrow{\bm{G}} \thetab \xrightarrow{\mathrm{SMPL\text{-}X}}
\Sigma \xrightarrow{\mathrm{Kirchhoff}} \hbody \xrightarrow{}
\mathrm{SINR} \xrightarrow{\log_2(1 + \cdot)} R
\end{equation}
is a composition of differentiable maps. The Kirchhoff render
$\Sigma \mapsto \hbody$ is linear in the per-triangle vertex
positions, so its Jacobian is trivial. SMPL-X is a sequence of
linear-blend-skinning matrix products with closed-form derivatives,
and the VPoser decoder is a feed-forward MLP with a continuous
6-D rotation representation. The gradient $\nabla_{\bm{z}} R$ is
computed end-to-end by reverse-mode autodifferentiation, with one
forward pass and one backward pass per iteration and no
finite-difference oracle on the underlying physics.

We identify the comfort region with the Mahalanobis ball
$\|\bm{z}\| \leq \rho_\mathrm{nat}$ in VPoser latent
space~\cite{Pavlakos2019VPoser}. This is an empirical surrogate for
the natural-pose manifold under the VAE prior. Calibration of
$\rho_\mathrm{nat}$ against subjective-comfort studies is left for
future work. The constraint defines
$C(\bm{z}) = \|\bm{z}\|^{2} \leq \rho_\mathrm{nat}^{2}$, anchored
at the VPoser prior origin and not at the user's current $\bm{z}_0$.
Under the VAE prior
$p(\bm{z}) = \mathcal{N}(\bm{0},\bm{I})$, the constraint
$\|\bm{z}\| \leq \rho_\mathrm{nat}$ keeps the prior log-probability
of every admissible pose above
$\log p(\bm{0}) - \rho_\mathrm{nat}^{2}/2$. The radius is set to
$\rho_\mathrm{nat} = 2.0$, which keeps poses within $2$~nats of the
prior mode and excludes the contortionist tail such as yoga splits
and gymnast bridges that VPoser assigns small density. The current
pose $\bm{z}_0$ enters only as the SINR baseline against which the
rate gain is measured.

\subsection{Human-in-the-loop control law}

The latent gradient $\nabla_{\bm{z}} R$ is computed on the device
and not actuated mechanically. The user is the actuator. A one-shot
pose target $\thetab^{*}$ would require the BS to know the user's
exact joint capabilities and current intent. The BS instead issues
at each tick a small latent step $\Delta\bm{z}$ along
$\nabla_{\bm{z}} R$, projected onto the naturalness ball
$\|\bm{z}\| \leq \rho_\mathrm{nat}$. The SMPL-X decoder maps
$\Delta\bm{z}$ to the corresponding skeletal configuration. The
user accepts or rejects, and the next tick observes the new IMU
pose, re-encodes it to a new $\bm{z}_0$, runs a fresh $\gammab$ fit
on the new uplink pilot, and emits the next step
(\Cref{alg:closed-loop}). The loop is projected gradient ascent on
a slowly drifting non-stationary objective with measurement
validation at each iterate.

\begin{algorithm}[!t]
\caption{$K$-iteration latent projected gradient ascent.}
\label{alg:closed-loop}
\begin{algorithmic}[1]
\State \textbf{Input:} baseline pose $\thetab^{(0)}$, naturalness
radius $\rho_\mathrm{nat}$, step size $\eta$, iteration budget $K$
\State Encode $\bm{z}_0 \leftarrow \bm{G}^{-1}(\thetab^{(0)})$
via the VPoser encoder
\State Initialize $\bm{z}^{(0)} = \bm{z}_0$ if $\|\bm{z}_0\| \leq \rho_\mathrm{nat}$ else
$\bm{z}^{(0)} = \bm{z}_0\,\rho_\mathrm{nat}/\|\bm{z}_0\|$
\For{$k = 0, 1, \ldots, K-1$}
\State Decode $\thetab^{(k)} \leftarrow \bm{G}(\bm{z}^{(k)})$ and run
SMPL-X forward map $\thetab^{(k)} \mapsto \Sigma(\thetab^{(k)})$
\State Render $\hbody^{(k)} \leftarrow$ UE-anchored Kirchhoff
integral on $\Sigma(\thetab^{(k)})$
\State Observe UL pilot residual $\Delta^{(k)} = \hh^{\mathrm{meas}}
- (\hBS + \hbody^{(k)})$
\State Fit $\gammab^{(k)}$ over the $K_{99}$-mode SVD basis of
$\Lambda_{\mathrm{KH}}^{(k)}$
\State Update predicted channel
$\hat\hh^{(k)} = \hBS + \hbody^{(k)} +
\Lambda_{\mathrm{KH}}^{(k)} \bm{V}_{K_{99}} \gammab^{(k)}$
\State Compute latent gradient $\bm{g}^{(k)} = \nabla_{\bm{z}}
R(\bm{z}^{(k)})$ by reverse-mode autodifferentiation through~\eqref{eq:autodiff-composition}
\State Project step onto the naturalness ball:
$\Delta\bm{z}^{(k)} = \mathrm{Proj}_{\|\bm{z}\|\leq \rho_\mathrm{nat}}\!
\bigl(\bm{z}^{(k)} + \eta\,\bm{g}^{(k)}\bigr) - \bm{z}^{(k)}$
\State User actuates: $\bm{z}^{(k+1)} = \bm{z}^{(k)} + \Delta\bm{z}^{(k)}$
\EndFor
\State \Return $\thetab^{(K)} = \bm{G}(\bm{z}^{(K)})$,
$R^{(K)} = R(\bm{z}^{(K)})$
\end{algorithmic}
\end{algorithm}

The latent gradient need not be exposed to the user. A one-bit
acceptance channel suffices to close the loop, with
$\Delta\bm{z}^{(k)}$ refreshed at $\sim 10$~Hz under body-scale
actuation. The control loop emits two semantic outputs at the user
side, namely the latent step $\Delta\bm{z}^{(k)}$
of~\Cref{alg:closed-loop} and the live exposure reading
of~\Cref{sec:saridx}.


% =====================================================================
\section{Numerical study}\label{sec:numerical}
% =====================================================================

\subsection{Existence, reachability, failure modes}

The closed loop of~\Cref{sec:closed-loop,sec:control} either reaches
a useful operating point or it does not. Three nested questions
structure the rest of this section.
\begin{enumerate}
\item[\textbf{(E)}] \textbf{Existence.} For a given (BS-geometry,
scene, baseline-pose) tuple, does any comfortable latent move
$\bm{z}$ with $\|\bm{z}\| \leq \rho_\mathrm{nat}$ produce a rate
gain $R(\bm{z}) - R(\bm{z}_0) \geq \Delta R_{\mathrm{tgt}}$?
\item[\textbf{(R)}] \textbf{Reachability.} Given that such a move
exists, does $K$-iteration UL-pilot-corrected projected gradient
ascent on $\bm{z}$ from the baseline reach within $1\,$dB of the
comfortable optimum?
\item[\textbf{(F)}] \textbf{Failure modes.} For trajectories that
fail to reach, what fails: stagnant ascent, a local-minimum trap,
calibration drift, or IMU noise?
\end{enumerate}

\subsection{Setup}

All experiments use the SMPL-X neutral mesh. Baseline poses are
sampled from $50$ AMASS walks at $30\,$fps, yielding
${\sim}13{,}000$ plausible pose frames. Each baseline frame
$\thetab^{0}$ is encoded to a VPoser latent $\bm{z}_0$ that
serves as the SINR baseline, not as a constraint anchor. The
naturalness ball $\|\bm{z}\| \leq \rho_\mathrm{nat} = 2.0$ keeps
every admissible pose within $2$ nats of the VPoser prior mode and
rules out the contortionist tail. The corpus is restricted to
walking sequences. It captures locomotion-pose statistics, not
deliberate pose changes such as raising an arm to obstruct or
unblock the link. Extension to gestural and body re-orientation
motions is left to future work.

The base station is an $8\times 8$ URA at $28\,$GHz with
$\lambda/2$ element spacing, mounted on a $35\,$m mast at
$(8.5, 21.7, 35)\,$m, transmitting $43\,$dBm against a noise
floor of $-94\,$dBm. The panel is downtilted and yawed to
broadside the served plaza, with panel-normal vector
$\hat{\bm{n}}_\mathrm{BS} = (0.30, 0.86, -0.41)$ pointing at the
ground centroid $(33, 91, 1.5)\,$m of the served LOS region. The
user body centroid is at $1.2\,$m height. The phone is held at
chest height, $40\,$cm in front of the body centroid toward the
BS. Rates are Shannon $\log_2(1 + \mathrm{SINR})$, with no MCS
cap.

\paragraph*{On-device pose estimation.} The on-device twin
consumes a per-tick SMPL-X pose estimate $\hat{\bm{\theta}}$
produced by a sparse-IMU body-pose inference pipeline running
on the user's consumer devices. Recent works in this lineage,
including DIP~\cite{DIP2018}, TransPose~\cite{TransPose2021},
IMUPoser~\cite{Mollyn2023IMUPoser}, and
MobilePoser~\cite{Xu2024MobilePoser}, regress full-body pose from
sparse inertial readings at real-time rates on a phone. The
reported consumer-only accuracy is ${\sim}24^{\circ}$ mean
per-joint rotation error (MPJRE) at $60\,$fps on a current
smartphone~\cite{Xu2024MobilePoser}. Research-grade body-suit
estimators (e.g., PIP at $6$ IMUs) reach ${\sim}12^{\circ}$. The
IMU sensitivity sweep (SI~\S{}S6) measures the closed-loop SINR
cost of operating on $\hat{\bm{\theta}}$ rather than the
noiseless $\bm{\theta}$ across per-joint attitude noise
$\sigma_\mathrm{joint} \in \{0,2,4,8,16\}^{\circ}$, spanning
research-grade body suits ($\sigma \approx 4^{\circ}$) through
current consumer-only operating points
($\sigma \approx 16^{\circ}$). The peak-APD location sensitivity
sweep (SI~\S{}S7) does the same for the cohort-side absorbed-power
reading.

The propagation environment is the canonical Munich scene
distributed with Sionna~RT~v$2.0.1$~\cite{Hoydis2023Sionna}, a
city-block model widely used as a benchmark for ray-tracing
studies.
We manually selected 18~Rx positions
(9~LOS, 9~NLOS, \Cref{fig:munich-scene}) from a 3D rendering
of the Sionna geometry. LOS positions sit on the open plaza
visible from the panel broadside. NLOS positions sit on side
streets shadowed by intervening buildings. The Rx span
$x \in [-69, 136]\,$m and $y \in [-27, 159]\,$m around the BS at
$(8.5, 21.7, 35)\,$m. Each Rx is then traced with Sionna~RT at
the carrier frequency $28\,$GHz, with maximum interaction depth
$5$, specular reflection, refraction, and edge diffraction
enabled (diffuse reflection disabled), $10^{7}$ samples per
source, and a per-source path-count cap of $2\times 10^{5}$.
Building and ground materials default to the Sionna scene library
values (concrete and the ITU \emph{itu\_concrete} variants). The
chosen LOS cells return ${\sim}500$ paths and the NLOS cells
${\sim}300$. The environmental twin does not need exact path gains.
The BS-side $\bm{\beta}$ calibration of~\Cref{sec:closed-loop}
absorbs the per-path amplitude and
phase discrepancies between the simulated scene and the live
channel, provided the ray-traced angles of arrival are
approximately right. The BS needs the angle of arrival of the
dominant MPCs at the served UE,
derivable from the UE position (GPS or IMU) and the scene
geometry. Per-path amplitudes and phases are calibrated
downstream. Reflection depth and edge diffraction are kept
enabled because deep-shadow NLOS paths reach the UE only after
several interactions. Per-element steering across the URA is
applied analytically via planewave phase ramps in the per-path
direction of departure from the array centroid, and the
body-mediated channel $\hbody$ is computed by feeding the
per-path body-incident bundle to the Kirchhoff render
of~\Cref{sec:kirchhoff}.

\begin{figure}[!t]
\centering
\includegraphics[width=\columnwidth]{rx_grid_annotated.png}
\caption{Canonical Munich scene from Sionna~RT, with the BS at
$(8.5, 21.7, 35)\,$m. The yellow arrow shows the panel broadside
direction. The $9$ LOS Rx (green) sit on the open plaza and side
avenues visible from the panel. The $9$ NLOS Rx (red) sit on
streets shadowed by intervening buildings. Rx were selected by
hand from a 3D rendering of the scene. Phone height $1.5\,$m;
body centroid $0.40\,$m behind the phone, facing the BS.}
\label{fig:munich-scene}
\end{figure}

\subsection{Deep-shadow example}

The deep-shadow example is shown in~\Cref{fig:deep-shadow-pose} on the
highest-gain Munich NLOS scene (scene~$17$,
$\mathrm{Rx} = (-69, 115, 1.5)\,$m on a deeply shadowed alley
north-west of the BS, scene loss $-109\,$dB). A latent
perturbation within the naturalness ball
$\|\bm{z}\|\leq\rho_\mathrm{nat}=2.0$ produces an $8.6\,$dB
cascaded-SINR difference between the encoder-projected baseline
pose ($10.9\,$dB) and the best of $N=100$ random latent samples
($19.5\,$dB). The visual difference between the two silhouettes
in~(a) and~(b) is small: a torso rotation and arm
shift of the order of $10$--$20^{\circ}$ at the dominant joints.
The whole-body azimuth is fixed by the per-Rx geometry (the body
faces the BS) and is not a control variable in this study. With
the BS-side bundle starved by the deep shadow, the body is the
dominant scatterer and a small change in its surface orientation
swings the cascaded SINR by $8.6\,$dB. Panels~(a) and~(b) show
the body silhouettes from the BS perspective with per-triangle
BS-illumination shading. Panel~(c) plots the SINR of each random
latent sample, with the baseline and best marked.

\begin{figure*}[!t]
\centering
\includegraphics[width=0.95\textwidth]{fig_hero_munich.pdf}
\caption{Deep-shadow example on Munich NLOS scene~$17$ (Rx at
$(-69, 115)\,$m on a deeply shadowed alley, scene loss
$-109\,$dB). (a)~Baseline pose (the VPoser encoder mean of an
AMASS walk frame), cascaded SINR $10.9\,$dB. (b)~Best of
$N=100$ random latent samples within the naturalness ball,
$19.5\,$dB. Each visible triangle is shaded by its
BS-illumination cosine. (c)~Cascaded SINR for all $100$ random
latent samples, with the baseline (red dashed) and best-found
(green dashed) marked at $\Delta = 8.6\,$dB.}
\label{fig:deep-shadow-pose}
\end{figure*}

\subsection{Existence}

\Cref{fig:existence} reports the per-scene SINR gain achieved by
a naturalness-bounded latent perturbation over the 18~Munich
scenes of~\Cref{fig:munich-scene}, stratified by LOS and NLOS.
The nine LOS scenes carry 53--58~dB of SINR from the BS multipath
bundle and admit small uplifts: gains span 0.6--2.4~dB
(mean 1.7~dB). The nine NLOS scenes split into three
regimes by scene loss. Mild-shadow NLOS (${\sim}75$ to $90\,$dB
scene loss) returns the same $1$ to $4\,$dB pose uplift as LOS.
Strong-shadow NLOS (${\sim}90$ to $100\,$dB) opens $4$ to $8\,$dB
of pose gain, where the body starts to contribute a
non-trivial scattering term. The deep-shadow case
(scene~$17$, Rx at $(-69, 115)\,$m, scene loss $-109\,$dB) is
the body-scatterer regime: the BS-side bundle alone delivers 11~dB
of SINR, and a naturalness-bounded pose move adds 8.6~dB, with
the body acting as the dominant scatterer back to the phone.

\begin{figure}[!t]
\centering
\includegraphics[width=\columnwidth]{fig_munich_perscene.pdf}
\caption{Per-scene cascaded SINR over the 18~Munich scenes
(9~LOS, 9~NLOS). Light bar: baseline pose $\bm{z}_0$
(VPoser-encoded AMASS walk frame). Dark bar: best of $N=100$
random latent samples within the naturalness ball
$\|\bm{z}\| \leq \rho_\mathrm{nat} =
2.0$~\cite{Pavlakos2019VPoser}. The deep-shadow scene~$17$
(rightmost NLOS) shows the body-scatterer regime, with +8.6~dB of
pose gain where the BS-side bundle alone is starved.}
\label{fig:existence}
\end{figure}

\subsection{Reachability}

The reachability of the projected latent gradient-ascent loop
of~\Cref{alg:closed-loop} at
$K \in \{2, 5, 10, 15, 20\}$ iterations is shown
in~\Cref{fig:reachability}. Success is defined as terminal SINR
within $1\,$dB of the brute-force optimum, taken as the
naturalness-ball maximum found by $N=100$ random samples. At
$K = 20$, the loop reaches the optimum on $8/9$ LOS scenes and
$8/9$ NLOS scenes. The two failures are dominated by local
plateaus in the latent-space rate landscape. In the deep-shadow
scene~$17$, the loop terminates within $1.1\,$dB of the
brute-force ball optimum. The IMU-noise sensitivity sweep
(SI~\S{}S6) shows that the rate loss the loop pays under
consumer-grade IMU operating points is small.

\begin{figure}[!t]
\centering
\includegraphics[width=\columnwidth]{fig_munich_reachability.pdf}
\caption{Reachability fraction over $18$ Munich scenes. The
closed loop of~\Cref{alg:closed-loop} is run for
$K \in \{2, 5, 10, 15, 20\}$ iterations of autodiff projected
gradient ascent on the VPoser latent. Success is terminal SINR
within $1\,$dB of the brute-force naturalness-ball optimum.}
\label{fig:reachability}
\end{figure}

\subsection{Per-joint sensitivity}

The joint-space sensitivity
$|\partial R / \partial \theta_j|$ across the $22$ SMPL-X body
joints, averaged over $50$ AMASS walk poses at the same canonical
scene, is shown in~\Cref{fig:per-joint}. The closed loop
optimizes in latent space, not joint space, but the joint-space
ranking is interpretable feedback for the user. The top five
joints are the pelvis, the two hips, and the two lower spine
joints. Together they carry over 80\% of the gradient mass.
The wrists, elbows, ankles, and head are flat: they do not move
the body-to-phone Kirchhoff coupling under normal phone-holding
geometries. The VPoser decoder respects this anatomy
automatically, so the latent gradient $\nabla_{\bm{z}} R$
concentrates on modes that actuate torso joints rather than
periphery joints.

\begin{figure}[!t]
\centering
\includegraphics[width=\columnwidth]{fig_per_joint_sensitivity.pdf}
\caption{Per-joint pose-rate sensitivity $|\partial R / \partial
\theta_j|$ ranked across the $22$ SMPL-X body joints, mean
$\pm$IQR over $50$ AMASS walk poses. The top five joints (pelvis,
left/right hip, spine1, spine2) carry the bulk of the gradient
mass. The anatomical inset shows their location on the body.
Periphery joints (wrists, elbows, ankles, head) are flat.}
\label{fig:per-joint}
\end{figure}

\subsection{Comfort-vs-rate Pareto}

The SINR gain achievable as a function of the latent comfort
radius $\|\bm{z} - \bm{z}_0\|$ is shown in~\Cref{fig:pareto},
sampled with $500$ random latent perturbations across the ball
$\|\bm{z} - \bm{z}_0\| \leq \rho_z = 1.5$ at one representative
NLOS-class scene from the parametric scalar-loss companion sweep
(SI~\S{}S3). The Pareto envelope is monotone in
$\|\bm{z}-\bm{z}_0\|$: a $0.5$-radius latent step buys
${\sim}6\,$dB of SINR gain, a unit step buys ${\sim}7.5\,$dB,
and the marginal return saturates above
$\|\bm{z}-\bm{z}_0\| \approx 1.2$. The cloud of below-frontier
samples shows that random comfort-bounded latent moves can also
lose SINR. The directed projected gradient ascent
of~\Cref{alg:closed-loop} tracks the frontier.

\begin{figure}[!t]
\centering
\includegraphics[width=\columnwidth]{fig_comfort_pareto.pdf}
\caption{Comfort-vs-rate Pareto envelope at one representative
NLOS-class scene (parametric scalar-loss companion sweep,
SI~\S{}S3). Random latent perturbations ($N=500$) fill the
cloud. The Pareto frontier (red) is the upper envelope. A
unit-radius latent step buys ${\sim}7.5\,$dB of SINR gain.}
\label{fig:pareto}
\end{figure}

\subsection{Comparison against baselines}

Four precoder families are compared in~\Cref{fig:baselines}
across the three scalar-loss bands of the parametric companion
sweep (SI~\S{}S3): no-twin zero-forcing (ZF), no-twin
maximum-ratio transmission (MRT), a T-pose-prior twin (the BS
assumes the canonical SMPL-X T-pose), and the full closed-loop body
twin. In the LOS-attenuated band, all twin-aware methods reach the
same operating point, and only no-twin ZF degrades because it inverts
against the wrong channel ($-18\%$ relative to no-twin MRT). In the
body-comparable band ($55\,$dB scalar loss), the proposed closed loop
reaches $7.1\%$ over no-twin MRT ($48\,$Mbps), and the T-pose-prior
twin reaches $5.8\%$. In the body-dominant band ($80\,$dB scalar
loss), both the T-pose twin and the proposed loop rescue the link
($+13.8\%$ over no-twin MRT). The
body-comparable band is where the body-side calibration most
clearly separates the body twin from a static-pose prior.

\begin{figure}[!t]
\centering
\includegraphics[width=\columnwidth]{fig_baselines.pdf}
\caption{Mean rate by scalar-loss band and precoder family on $30$
AMASS scenes ($10$ per band) from the parametric companion sweep
(SI~\S{}S3). Error bars are $\pm$one standard deviation
across scenes. No-twin ZF degrades in the LOS-attenuated band.
The proposed closed loop separates from the static T-pose prior in
the body-comparable band, where the body-side calibration is
decisive.}
\label{fig:baselines}
\end{figure}


% =====================================================================
\section{Live exposure reading and RF-EMF cohort studies}\label{sec:saridx}
% =====================================================================

\subsection{Absorbed-power reading}

At each twin tick, the on-device twin evaluates~\Cref{eq:Sab-norm}
on the visible triangles for the current precoder $\xx$ and
integrates over the body surface to obtain the whole-body absorbed
power $\Pabs$ and the cumulative dose $\int\Pabs\,dt$ over the session.
The reading is displayed in milliwatts (instantaneous) and
milli-watt-hours (cumulative). At the canonical operating point
($43\,$dBm BS power, $30\,$m range, $28\,$GHz,
$8\times 8$~URA), $\Pabs$ is $0.58\%$ of the ICNIRP
general-public reference of $4\,$W/m$^{2}$ surface-equivalent
power density~\cite{ICNIRP2020}. The link budget, not compliance,
therefore binds $\xx$ at the FR2 operating points considered here.
The reading still provides a per-slot exposure value that current
regulatory and epidemiological workflows do not measure directly.

The reading reported here is the whole-body integrated $\Pabs$
and its time integral. Both are robust to on-device pose error
(SI~\S{}S7). At the consumer-grade IMU operating point
($\sigma_\mathrm{joint} = 16^{\circ}$), the on-device $\Pabs$
tracks the true value to within a $-3.2\%$ mean bias and
$\pm 4.5\%$ standard deviation across the 18~Munich scenes;
at $\sigma_\mathrm{joint} = 4^{\circ}$ (a research-grade
6-IMU body suit), the bias drops to $-0.5\%$ with
$\pm 1.4\%$ standard deviation. The peak-local $\Sab$ over the ICNIRP
$4\,$cm$^{2}$ surface-average window is \emph{not} reported
per-tick. The argmax over body triangles is fragile under
realistic IMU pose noise: the peak triangle centroid moves by a
mean of 52--58~cm across $\sigma_\mathrm{joint} \in
\{2,16\}^{\circ}$ (SI~\S{}S7), so per-tick anatomical
localization of the hot spot is out of scope. Time-aggregated or
larger-window peak-$\Sab$ statistics may recover a stable
localized reading. We leave that to follow-up work.

\subsection{One operator, two semantic outputs}

The reading is computed on the same per-triangle Kirchhoff
operator as $\hbody$, with the integrand evaluated in absorption
mode ($\Sab = \|\tilde{\bm{G}}\xx\|^{2}$) rather than radiation
mode ($\hbody = \int \tilde{\bm{G}}\xx \cdot e^{-jkr}/r\,\mathrm{d}A$).
Both modes share the operator factorization~\Cref{eq:Q-form} and
the same per-element exposure channel matrix $\tilde{\bm{G}}$.
The SAR-matrix method~\cite{Hochwald2014,Ying2015,Ying2017}
that supplies $\Qabs$ to the BS-side exposure-constrained
beamformer is the same method that supplies $\Sab(\rr)$ to
the user-facing reading. The reading is one read-out mode of an
object the network already maintains for precoding.

\subsection{Aggregation for cohort exposure studies}

Prospective RF-EMF cohort studies link bio/health outcomes in
well-characterized populations to RF-EMF
exposure~\cite{Roosli2010,Roosli2021,vanWel2024}. The exposure
side currently rests on self-reported phone-use questionnaires
and modeled SAR proxies. The European CLUE-H
consortium~\cite{CLUEH2025} extends the single-cohort lineage of
COSMOS~\cite{Roosli2021} and HERMES~\cite{Roosli2010} to a
coordinated set of EU cohorts and binds on this same exposure
side: questionnaire-reported phone use and modeled SAR proxies
are coarse compared to the bio-outcome resolution the cohorts
target.

Per-individual absorbed-energy time series, opt-in-aggregated
from the reading, upgrade the exposure side from self-report to
direct measurement, with no in-cohort instrumentation. The
privacy delta is small: only the scalar $\Pabs$ time series need
leave the device for cohort use, with no body-twin parameters
and no surface deposition maps. Aggregation cadence is the
user's choice. Longitudinal cohort studies want monthly or
quarterly summaries, not per-slot streaming.
\Cref{fig:sab-reading} illustrates the on-device computation:
a per-triangle $\Sab$ map on the SMPL-X body, an instantaneous
$\Pabs$ time series across a $30\,$s walk, and the $100$-user
cumulative-dose distribution at the canonical operating point.

\begin{figure*}[!t]
\centering
\includegraphics[width=0.95\textwidth]{fig_sab_receipt.pdf}
\caption{Live exposure reading on the SMPL-X body twin.
(a)~Per-triangle $\Sab$ heatmap on the body surface for one
representative pose. The cyan star marks the phone position.
(b)~Instantaneous $\Pabs$ over a $30\,$s AMASS-walk trace, with
cumulative dose on the secondary axis. (c)~ECDF of $30\,$s
cumulative dose across $100$ random users. The orange dashed
line marks the $4\,$W/m$^{2}$ ICNIRP general-public
reference~\cite{ICNIRP2020}. The user distribution lies well
below it.}
\label{fig:sab-reading}
\end{figure*}


% =====================================================================
\section{Discussion}\label{sec:discussion}
% =====================================================================

\subsection{Bandwidth and per-slot compute}\label{sec:bandwidth}

\Cref{tab:cadence} lists the refresh cadences and per-update
wall-clock times of the on-device twin components. The
UE-anchored Kirchhoff render takes $\sim\!100\,$ms per pose on a
single CPU core at $T_{\mathrm{vis}} \sim\!10^{4}$ visible
triangles. The body-side $\bm{\gamma}$ calibration is a ridge
least-squares fit over the $K_{99}$-mode SVD basis and takes
$\sim\!1\,$ms. The pose-gradient finite difference over the top
six sensitivity-ranked joints requires seven Kirchhoff renders,
$\sim\!700\,$ms total. The downlink pose-suggestion payload is
$K\log_2(M)$ bits per twin tick, which at typical $K\leq 4$ and
$M=64$ amounts to $\sim\!24\,$bits and sits well below the
5G~NR downlink control-channel budget.

\begin{table}[!t]
\caption{Refresh cadences and per-update wall-clock for the
on-device twin loop.}
\label{tab:cadence}
\centering\scriptsize
\begin{tabular}{@{}p{0.47\columnwidth} c p{0.24\columnwidth}@{}}
\toprule
Component & Cadence & Per-update wall-clock \\
\midrule
SMPL-X forward map $\thetab \mapsto \Sigma$ & $\sim\!10\,$Hz & $\sim\!5\,$ms \\
UE-anchored Kirchhoff render & $\sim\!10\,$Hz & $\sim\!100\,$ms \\
$\bm{\gamma}$ ridge LS over $K_{99}$ basis & $\sim\!1\,$Hz & $\sim\!1\,$ms \\
Pose gradient (FD on top-6 joints) & $\sim\!1\,$Hz & $\sim\!700\,$ms \\
DL pose-suggestion payload & per-tick & $K \log_2(M) \approx 24\,$bits \\
UL pilot for residual & per-slot & embedded in CSI feedback \\
\bottomrule
\end{tabular}
\end{table}

\subsection{Privacy delta relative to a model-only baseline}

The SMPL-X pose vector $\thetab$ does not leave the device. The
BS sees only the cascaded channel $\hh^{\mathrm{meas}}$ via
uplink pilots and the body-side calibration coefficient
$\gammab\in\Complex^{K}$ for $K\leq M$. The pose suggestion the
BS issues is a direction in pose space, not a target pose. The
per-tick $\Pabs$ scalar is the only body-side telemetry that
leaves the device, and only on opt-in for cohort use.

\subsection{Limitations}

First, the analysis is single-user. Multi-user coexistence, in
which one user's pose move redirects another user's multipath,
is not treated here.
Second, the back-face cull suffices for torso, head, and
limb-scale features at FR2. Fingertip features and ear edges
require strict ray-mesh visibility, which raises the per-pose
render cost by a factor of $\sim\!5$ and is detailed in the
supplementary information.
Third, the sub-6~GHz extension via the layered Fabry--P\'erot
transmission model introduces frequency-selective effects, in
particular the $\lambda/4$ fat-resonance dip, that the mmWave PO
regime does not exhibit. The supplementary information treats
this case explicitly.

\subsection{Future work}

Three extensions follow directly from these limitations. The
single-user restriction lifts once the cascaded channel
carries multiple body twins, in which case one user's pose move
becomes a controlled multipath perturbation for another user.
Replacing the SMPL-X mesh with a chassis parametric model
extends the same Kirchhoff render to vehicular twins. The static
MCS cap used here can be replaced by an application-side rate
target, so that the pose suggestion adapts to the AR/XR frame
budget rather than to a fixed throughput floor.


% =====================================================================
\section{Conclusion}\label{sec:conclusion}
% =====================================================================

We presented an on-device digital twin of the user's body
embedded in the per-slot control loop of an FR2 radio access
network. From the same per-triangle Kirchhoff--physical-optics
machinery the twin produces a live exposure reading and a
pose-differentiable mmWave channel, and uplink-pilot calibration
through a $K$-mode SVD basis closes the loop in both directions.
A UE-anchored Kirchhoff render evaluates the body-mediated
channel coefficient in $\mathcal{O}(M\,T_{\mathrm{vis}})$
operations per pose and is the receiver-anchored dual of forward
radiosity at small receiver counts. The on-device loop runs at
$\sim\!1$ to $10\,$Hz on a single CPU core.

Two consequences follow for FR2 deployment. The cascaded-SINR
gain from a naturalness-bounded pose move is large enough on
deep-shadow NLOS scenes to convert outage to service, and the
peak-local absorbed power density at the canonical operating
point sits two orders of magnitude below the ICNIRP
general-public reference. The body twin therefore returns the
per-individual exposure time series that the existing
regulatory framework specifies but does not measure directly,
the quantity that prospective RF-EMF cohort studies otherwise
approximate via self-report.

The single-user restriction is the next item to lift, with the
same Kirchhoff render carrying multiple twins on the cascaded
channel. The same Kirchhoff render transfers to vehicular
twins under a chassis parametric model, and to AR/XR
applications by replacing the static MCS cap with an
application-side rate target.


% =====================================================================
\appendices
% =====================================================================
\section{Per-path Fresnel transmission operator}\label{app:fresnel}

This appendix gives the formal definition of the Fresnel transmission
operator $\Fn$ that appears throughout
\Cref{sec:physics,sec:kirchhoff}. The derivation from Maxwell's
boundary conditions follows the standard route~\cite{BornWolf}.

At the surface point $\bm{r}$ on triangle $t$, path $n$ has
incidence cosine $\mu_n = \hat{\bm{n}}_t \cdot (-\hat{\bm{k}}_n)$.
For a front-facing path ($\mu_n > 0$) the TE and TM unit vectors are
\begin{equation}\label{eq:TE-TM-basis}
  \hat{\bm{e}}_{s,n} = \frac{\hat{\bm{k}}_n \times \hat{\bm{n}}_t}
                            {\|\hat{\bm{k}}_n \times \hat{\bm{n}}_t\|},
  \qquad
  \hat{\bm{e}}_{p,n} = \hat{\bm{e}}_{s,n} \times \hat{\bm{k}}_n\,,
\end{equation}
and the Fresnel amplitude transmission coefficients are
\begin{equation}\label{eq:fresnel-t}
  t_{s,n} = \frac{2\mu_n}{\mu_n + \xi_n}\,, \qquad
  t_{p,n} = \frac{2\,\ntilde\,\mu_n}{\ntilde^2\,\mu_n + \xi_n}\,,
\end{equation}
with $\xi_n \equiv \sqrt{\ntilde^2 - 1 + \mu_n^2}$ and
$\RE(\xi_n) > 0$. For tissue at 28~GHz, the IT'IS database gives
$\ntilde \approx 4.49 - 1.79j$ for skin.

\begin{definition}[Fresnel transmission operator]\label{def:fresnel-op}
The dimensionless rank-2 operator
\begin{equation}\label{eq:Fn-def}
  \Fn(\bm{r}) =
  \bigl(t_{s,n}\,\hat{\bm{e}}_{s,n}\hat{\bm{e}}_{s,n}^{\!\top}
       + t_{p,n}\,\hat{\bm{e}}_{p,n}\hat{\bm{e}}_{p,n}^{\!\top}\bigr)\,H(\mu_n)
\end{equation}
projects $\psin$ onto the local TE and TM directions and scales each
by its Fresnel transmission coefficient. The Heaviside factor
$H(\mu_n)$ vanishes on back-facing paths.
\end{definition}

The transmitted field at depth $z$ inside the tissue is then
\begin{equation}\label{eq:Etrans}
  \EE_n^{\mathrm{trans}}(\bm{r},z)
  = \Fn(\bm{r})\,\psin\,x_{j(n)}\,
    e^{-jk_0\hat{\bm{k}}_n\cdot\bm{r}}\,
    e^{-jk_0\xi_n z}\,,
\end{equation}
where $x_{j(n)}$ is the precoder weight on the source antenna of
path $n$ and $j$ is the engineering imaginary unit. Writing
$k_0\xi_n = \beta_n - j\alpha_n$ with $\alpha_n > 0$ separates the
amplitude decay rate $\alpha_n$ (skin depth $\delta_n = 1/\alpha_n$)
from the in-tissue phase rate $\beta_n$. The per-triangle integrand
of the UE-anchored Kirchhoff render specialises~\eqref{eq:Etrans}
at $z = 0$ and contracts against the UE steering vector
$\bm{g}_{\mathrm{UE}}(\hat{\bm{\eta}})$.

The operator factorization~\eqref{eq:Fn-def} is the structural
upgrade over per-element scalar Fresnel calibrations (a single
$\beta\approx 0.70$ that lumps Fresnel transmission with depth
decay): it preserves the polarization rotation that the BS
precoder $\bm{x}$ can exploit at the per-triangle level, and it
isolates the frequency-dependent depth integral into the universal
coupling constant of~\Cref{prop:approx2-app}.

% =====================================================================
\section{Approximations 1 and 2 error budget}\label{app:approxs}

The reduction of the coherent absorption double sum to a single
squared norm rests on two approximations, restated and quantified
here. Full proofs are in SI~\S{}S1.

\begin{proposition}[Approximation 1: TM polarization direction]
\label{prop:approx1-app}
The exact transmitted TM field points along the refracted direction
$\hat{\bm{e}}'_{p,n}$ inside tissue, with
$\sin\theta_{t,n} = \sin\theta_n / |\ntilde|$ by Snell's law. The
Fresnel operator $\Fn$ in~\eqref{eq:Fn-def} replaces
$\hat{\bm{e}}'_{p,n}$ by the incident-side $\hat{\bm{e}}_{p,n}$.
Self-terms are unaffected
($\|\hat{\bm{e}}_{p,n}\|^2 = \|\hat{\bm{e}}'_{p,n}\|^2 = 1$).
The error on TM--TM cross-terms is $O(1/|\ntilde|^2)$, which is
$\lesssim 4\%$ for skin at 28~GHz.
\end{proposition}

The physical content is that the in-tissue normal field
$E_\perp \propto \sin\theta_{t,n} = O(1/|\ntilde|)$, so its
contribution to $\|\EE^{\mathrm{trans}}\|^2$ is $O(1/|\ntilde|^2)$.
The high tissue refractive index suppresses normal coupling, not
the polarization direction itself.

\begin{proposition}[Approximation 2: universal depth coupling]
\label{prop:approx2-app}
Define the normalized depth coupling
\begin{equation}\label{eq:Gamma-app}
  \Gamma_{nn'} \equiv
    \frac{2\sqrt{\alpha_n\alpha_{n'}}}
         {\alpha_n + \alpha_{n'} - j(\beta_{n'} - \beta_n)}\,,
\end{equation}
so that $\Gamma_{nn}=1$. Setting
$\Gamma_{nn'} \approx 1$ for all pairs reduces the coherent double
sum (SI~\S{}S1, eq.~(S2)) to a single squared norm,
and incurs a non-conservative error $\leq 0.5\%$ on every
cross-term for skin at 28~GHz.
\end{proposition}

The mechanism is that $|\ntilde| \approx 4.83$ for skin compresses
all refraction angles to within $\sim\!10^{\circ}$ of normal, so
$\alpha_n$ and $\beta_n$ are nearly angle-independent and
$\Gamma_{nn'}$ stays close to unity. The error is bounded by
sweeping all incidence-angle pairs.

The combined error budget across the millimeter-wave band is
summarized in~\Cref{tab:approx-acc}.

\begin{table}[!htbp]
\centering
\caption{Combined error budget for Approximations 1 and 2 on
biological skin, from the IT'IS v5.0 database. App.~1 error is
$\max_\theta \|\hat{\bm{e}}'_{p,n} - \hat{\bm{e}}_{p,n}\|$ projected
onto cross-terms. App.~2 error is $\max_{n,n'}|1 - \Gamma_{nn'}|$.
Computed from the analysis of SI~\S{}S1.}
\label{tab:approx-acc}
\begin{tabular}{rrrr}
\toprule
$f$ (GHz) & $|\ntilde|$ & App.~1 (\%) & App.~2 (\%) \\
\midrule
10  & 5.87 & $\lesssim 3$  & $\leq 0.30$ \\
28  & 4.83 & $\lesssim 4$  & $\leq 0.44$ \\
40  & 4.30 & $\lesssim 5$  & $\leq 0.55$ \\
60  & 3.68 & $\lesssim 7$  & $\leq 0.75$ \\
100 & 3.01 & $\lesssim 11$ & $\leq 1.1$  \\
\bottomrule
\end{tabular}
\end{table}

Both errors increase with frequency as $|\ntilde|$ decreases, but
remain inside the $\sim\!20\%$ tissue-dielectric uncertainty across
the FR2 band. At 28~GHz they are an order of magnitude smaller than
the calibration headroom absorbed by the body-side $\gammab$
regression in~\Cref{sec:closed-loop}.

% =====================================================================
\section{Identifiability of \texorpdfstring{$K$}{K}-mode SVD calibration}
\label{app:ident}

This appendix gives the truncated-SVD argument behind~\Cref{prop:ident}.

Let $\Lambda_{\mathrm{KH}} \in \Complex^{M \times T_{\mathrm{vis}}}$
be the per-pilot Kirchhoff regressor matrix introduced in
\Cref{sec:closed-loop}, with thin SVD
$\Lambda_{\mathrm{KH}} = \bm{U}\bm{\Sigma}\bm{V}^{\!H}$ and singular
values $\sigma_1 \geq \sigma_2 \geq \cdots \geq \sigma_M > 0$. The
ridge least squares with body-side basis
$\bm{V}_K = \bm{V}_{:,1:K}$ solves
\begin{equation}\label{eq:gamma-ridge}
  \hat{\gammab} = \argmin_{\gammab \in \Complex^K}
    \bigl\|\bm{r} - \Lambda_{\mathrm{KH}}\bm{V}_K\gammab\bigr\|^2
  + \lambda \|\gammab\|^2\,,
\end{equation}
where $\bm{r} = \hat{\hh} - \hBS - \hbody \in \Complex^M$ is the
pilot residual after the LOS and nominal-pose body terms are
subtracted, and $\lambda \geq 0$ is a small ridge.

\textbf{Identifiability.} Substituting the SVD and using
$\bm{V}_K^{\!H}\bm{V}_K = \bm{I}_K$,
\begin{equation}\label{eq:Fb-Vk}
  \Lambda_{\mathrm{KH}}\bm{V}_K = \bm{U}_K\bm{\Sigma}_K\,,
\end{equation}
where $\bm{U}_K = \bm{U}_{:,1:K}$ and
$\bm{\Sigma}_K = \mathrm{diag}(\sigma_1,\dots,\sigma_K)$. The
unregularized normal equations give
$\hat{\gammab} = \bm{\Sigma}_K^{-1}\bm{U}_K^{\!H}\bm{r}$, which is
unique whenever $\sigma_K > 0$, i.e.~whenever the truncated basis
sits in the column space of $\Lambda_{\mathrm{KH}}$. If the residual
$\bm{r}$ lies entirely in $\mathrm{span}(\bm{U}_K)$, the recovery is
exact. Otherwise, the orthogonal complement
$(\bm{I}-\bm{U}_K\bm{U}_K^{\!H})\bm{r}$ lies outside
the body-side estimable subspace and is absorbed into the
LOS / antenna calibration $\betab$ in~\Cref{sec:closed-loop}.

\textbf{Noise amplification.} Under additive complex Gaussian pilot
noise $\bm{n} \sim \mathcal{CN}(\bm{0},\sigma_p^2 \bm{I}_M)$,
the per-mode estimation variance is
\begin{equation}\label{eq:gamma-var}
  \mathrm{Var}(\hat\gamma_k)
    = \frac{\sigma_p^2}{\sigma_k^2 + \lambda}\,,
\end{equation}
so the fit is well-posed in the per-mode SNR
sense exactly when $\sigma_k^2 \,\mathrm{SNR}_{\mathrm{ul}}\gg 1$.
Equivalently, the $K$-th mode is informative when the per-mode SNR
$\sigma_k^2/\sigma_p^2$ exceeds unity. This justifies the operating
choice $K = K_{99}$, defined as the largest $K$ with
$\sum_{k\leq K}\sigma_k^2 \geq 0.99\,\|\bm{\Sigma}\|_F^2$ and
clipped at $\sigma_K^2 / \sigma_p^2 \geq 1$. The $K_{99}$ basis is
adaptive to the BS--user--body geometry: rank-1 in the far-BS regime
where one body-bounce mode dominates, rank-many in the close-BS
regime where the body subtends a wide solid angle and several body
modes carry comparable spectral mass (operator-spectrum sweep is in
the SI).

\textbf{Lower bound on pilot SNR.} The fit is non-trivial when
$\mathrm{SNR}_{\mathrm{ul}} \geq K\sigma_K^{-2}$, i.e.~when the
per-mode noise variance~\eqref{eq:gamma-var} is small enough that
the $K$ leading modes are jointly resolvable. For the served-user
geometries swept in~\Cref{sec:numerical}, $K_{99} \in \{1,2,4\}$
and $\sigma_K^{-2} \leq 1$, so the operating point of
$20\,$dB UL pilot SNR sits inside the identifiable
region. \hfill$\square$


% =====================================================================
% Bibliography
% =====================================================================
\begin{thebibliography}{99}
\bibitem{MacCartney2017} G.~R. MacCartney \emph{et~al.}, ``Rapid
fading due to human blockage in pedestrian crowds at 5G
millimeter-wave frequencies,'' in \emph{Proc. IEEE Global Commun.
Conf. (GLOBECOM)}, Singapore, Dec.~2017, pp.~1--7.
\bibitem{Maccartney2017blockage} G.~R. MacCartney, T.~S. Rappaport,
and S.~Rangan, ``Path loss models for 5G millimeter wave
propagation channels in urban microcells,'' \emph{IEEE Access},
vol.~6, pp.~17\,460--17\,479, 2018.
\bibitem{ICNIRP2020} International Commission on Non-Ionizing
Radiation Protection, ``Guidelines for limiting exposure to
electromagnetic fields (100~kHz to 300~GHz),'' \emph{Health Phys.},
vol.~118, no.~5, pp.~483--524, May~2020.
\bibitem{SMPLX2019} G.~Pavlakos \emph{et~al.}, ``Expressive body
capture: 3D hands, face, and body from a single image,'' in
\emph{Proc. IEEE/CVF Conf. Comput. Vis. Pattern Recognit. (CVPR)},
Long Beach, CA, USA, Jun.~2019, pp.~10\,975--10\,985.
\bibitem{Pavlakos2019VPoser} N.~Ghorbani \emph{et~al.}, ``VPoser:
variational human pose prior for body inverse kinematics,''
software release, Max Planck Inst. Intell. Syst., T\"ubingen,
Germany, 2019. [Online]. Available:
\url{https://github.com/nghorbani/human_body_prior}
\bibitem{AMASS2019} N.~Mahmood \emph{et~al.}, ``AMASS: archive of
motion capture as surface shapes,'' in \emph{Proc. IEEE/CVF Int.
Conf. Comput. Vis. (ICCV)}, Seoul, South Korea, Oct.~2019,
pp.~5442--5451.
\bibitem{Hoydis2023Sionna} J.~Hoydis \emph{et~al.}, ``Sionna~RT:
differentiable ray tracing for radio propagation modelling,'' in
\emph{Proc. IEEE Globecom Workshops}, Kuala Lumpur, Malaysia,
Dec.~2023.
\bibitem{DiRenzo2020} M.~Di~Renzo \emph{et~al.}, ``Smart radio
environments empowered by reconfigurable intelligent surfaces:
how it works, state of research, and the road ahead,''
\emph{IEEE J.~Sel. Areas Commun.}, vol.~38, no.~11,
pp.~2450--2525, Nov.~2020.
\bibitem{Wu2019} Q.~Wu and R.~Zhang, ``Towards smart and
reconfigurable environment: intelligent reflecting surface aided
wireless network,'' \emph{IEEE Commun. Mag.}, vol.~58, no.~1,
pp.~106--112, Jan.~2020.
\bibitem{Madgwick2011} S.~O.~H. Madgwick, A.~J.~L. Harrison, and
R.~Vaidyanathan, ``Estimation of IMU and MARG orientation using a
gradient descent algorithm,'' in \emph{Proc. IEEE Int. Conf.
Rehabil. Robot. (ICORR)}, Zurich, Switzerland, Jun.--Jul.~2011,
pp.~1--7.
\bibitem{TS38214} \emph{NR; Physical Layer Procedures for Data},
3GPP TS~38.214, Rel.~17, Table~5.1.3.1-2, 2024.
\bibitem{Hochwald2014} B.~M. Hochwald, D.~J. Love, S.~Yan, P.~Fay,
and J.-M. Jin, ``Incorporating specific absorption rate constraints
into wireless signal design,'' \emph{IEEE Commun. Mag.}, vol.~52,
no.~9, pp.~126--133, Sep.~2014.
\bibitem{Ying2015} D.~Ying, D.~J. Love, and B.~M. Hochwald, ``Closed
loop precoding and capacity analysis for multiple-antenna wireless
systems with user radiation exposure constraints,'' \emph{IEEE
Trans. Wireless Commun.}, vol.~14, no.~10, pp.~5859--5870,
Oct.~2015.
\bibitem{Ying2017} D.~Ying, D.~J. Love, and B.~M. Hochwald,
``Beamformer optimization with a constraint on user electromagnetic
radiation exposure,'' \emph{IEEE Trans. Signal Process.}, vol.~65,
no.~9, pp.~2376--2389, May~2017.
\bibitem{Roosli2010} M.~R\"o\"osli, P.~Frei, E.~Mohler, and K.~Hug,
``Systematic review on the health effects of exposure to
radiofrequency electromagnetic fields from mobile phone base
stations,'' \emph{Bull. World Health Organ.}, vol.~88, no.~12,
pp.~887--896F, Dec.~2010.
\bibitem{Roosli2021} M.~R\"o\"osli \emph{et~al.}, ``Cosmos: cohort
study of mobile phone use and health,'' \emph{Environ. Int.},
vol.~149, art.~no.~106328, Apr.~2021.
\bibitem{BornWolf} M.~Born and E.~Wolf, \emph{Principles of
Optics}, 7th~ed.\quad Cambridge, U.K.: Cambridge Univ. Press,
1999.
\bibitem{Chen2024DTN} M.~Chen \emph{et~al.}, ``Digital twin
networks: a survey,'' \emph{IEEE Internet Things J.}, 2024.
\bibitem{Saad2025DTN} W.~Saad \emph{et~al.}, ``Digital twins for
6G wireless networks,'' \emph{Proc. IEEE}, 2025, to appear.
\bibitem{Hashash2025DTN} O.~Hashash \emph{et~al.}, ``Wireless
digital twins for cross-layer optimization,'' \emph{IEEE Commun.
Mag.}, 2025, to appear.
\bibitem{CLUEH2025} CLUE-H Consortium, ``Cohort linkage on
mobile-phone RF-EMF exposure and health outcomes,'' Horizon Europe
project, 2025. [Online]. Available: \url{https://clue-h.eu}
\bibitem{vanWel2024} L.~van~Wel \emph{et~al.}, ``Methods for
assessing personal RF-EMF exposure in prospective cohort studies,''
\emph{Environ. Int.}, vol.~189, art.~no.~108796, Jul.~2024.
\bibitem{general5G} \emph{Framework and Overall Objectives of the
Future Development of IMT for 2030 and Beyond}, ITU-R
Recommendation M.2160-0, Nov.~2023.
\bibitem{HeathLozanoFA2018} R.~W. Heath~Jr. and A.~Lozano,
\emph{Foundations of MIMO Communication}.\quad Cambridge, U.K.:
Cambridge Univ. Press, 2018.
\bibitem{BrusselsArrete2007} ``Arr\^et\'e de la R\'egion de
Bruxelles-Capitale relatif \`a la protection de l'environnement
contre les \'eventuels effets nocifs et nuisances provoqu\'es par
les radiations non ionisantes,'' Brussels, Belgium, Mar.~2007.
\bibitem{BrusselsArrete2024} ``Arr\^et\'e de la R\'egion de
Bruxelles-Capitale modifiant l'arr\^et\'e du 30~mars 2007,''
Brussels, Belgium, 2024.
\bibitem{ItalyDL} ``Decreto sui limiti di esposizione ai campi
elettromagnetici per le reti mobili,'' Gazzetta Ufficiale, Rome,
Italy, 2024.
\bibitem{DIP2018} Y.~Huang, M.~Kaufmann, E.~Aksan, M.~J. Black,
O.~Hilliges, and G.~Pons-Moll, ``Deep inertial poser: learning to
reconstruct human pose from sparse inertial measurements in real
time,'' \emph{ACM Trans. Graph.}, vol.~37, no.~6, art.~no.~185,
Dec.~2018.
\bibitem{TransPose2021} X.~Yi, Y.~Zhou, and F.~Xu, ``TransPose:
real-time 3D human translation and pose estimation with six
inertial sensors,'' \emph{ACM Trans. Graph.}, vol.~40, no.~4,
art.~no.~86, Aug.~2021.
\bibitem{Mollyn2023IMUPoser} V.~Mollyn, R.~Arakawa, M.~Goel,
C.~Harrison, and K.~Ahuja, ``IMUPoser: full-body pose estimation
using IMUs in phones, watches, and earbuds,'' in \emph{Proc. CHI
Conf. Hum. Factors Comput. Syst.}, Hamburg, Germany, Apr.~2023.
\bibitem{Xu2024MobilePoser} V.~Xu \emph{et~al.}, ``MobilePoser:
real-time full-body pose estimation and 3D human translation
from IMUs in mobile consumer devices,'' in \emph{Proc. ACM Symp.
User Interface Softw. Technol. (UIST)}, Pittsburgh, PA, USA,
Oct.~2024.
\bibitem{SwissNISV2000} ``Verordnung \"uber den Schutz vor
nichtionisierender Strahlung (NISV),'' Swiss Federal Council,
Bern, Switzerland, SR~814.710, Dec.~1999, with subsequent
amendments.
\bibitem{Christ2010VF} A.~Christ \emph{et~al.}, ``The Virtual
Family: development of surface-based anatomical models of two
adults and two children for dosimetric simulations,'' \emph{Phys.
Med. Biol.}, vol.~55, no.~2, pp.~N23--N38, Jan.~2010.
\end{thebibliography}

\end{document}
